\section{Conclusion}\label{sec:conclusion}

HOPPER makes a satoshi the protagonist of a game whose rules live on
Bitcoin. A mint inscription gives the sat a name and a face. A four-leaf
Taproot output gives it custody: three leaves sign a hop, one per direction,
and the fourth sends it onward. A 55-byte body binds the crossing to the sat,
the mint and the chosen direction, and the hash of block $T+1$ wraps to one
direction for the whole block. The v8 fold proves every body-valid crossing
against the raw transaction that spent the canonical leaf; when it cannot,
it says \emph{unverified} instead of guessing. The indexer and the wallet
make this usable, but neither is the referee.

The design is intentionally narrow. There is no token, no bridge, no
server-side game loop, and no privileged signer. The consensus surface is
small enough to specify in bytes and to test across two languages: a
generation sawtooth, a four-byte modulo, a direction comparison, and four
fixed leaves, backed by a proof check on every counted crossing. The art is a
self-contained document that can fail closed. The one unavoidable extra
transaction is a consequence of self-reference on a hash-linked ledger, not
an accident. The single grandfathered v7 Hopper is a custody exception
carved out of already-confirmed history; it does not widen the v8 rules.

The remaining work is engineering discipline: keep the byte fixtures
authoritative, keep the two implementations in lockstep, and keep the
indexer's conveniences out of the trust path. If those hold, a frog is what
the chain says it is. \textit{One satoshi. One life. The
chain is the referee.}
